% =====================================================================
\section{Development Environment and Technology Stack}
\label{sec:stack}
% =====================================================================

\subsection{Framework decisions}
The team will implement LipScribe in Python and select one tool for each role
to keep the codebase small and the learning burden manageable. Two choices
shape the rest of the project.

\paragraph{Deep-learning framework: PyTorch \rec.}
PyTorch provides every model component the proposal requires---3D
convolutions (\texttt{nn.Conv3d}), bidirectional LSTMs with packed
variable-length sequences, multi-head attention
(\texttt{nn.MultiheadAttention}) and a built-in CTC loss (\texttt{nn.CTCLoss}).
Its eager execution makes tensor shapes easy to inspect during debugging, and
\texttt{torchaudio} supplies log-Mel and MFCC transforms that run identically
on CPU and GPU. \texttt{torchvision} supplies an ImageNet-pretrained
EfficientNet-B0~\cite{tan2019} for the backbone experiment in \cref{sec:visual}.
TensorFlow could serve equally well; the team chooses PyTorch so that one
framework covers training, evaluation and inference.

\paragraph{Interface framework: Gradio \rec.}
Gradio builds a browser-based demonstration from Python functions, with
ready-made webcam, microphone and video-upload components. This matches the
``simple web demonstration'' promised in the proposal \spec\ without requiring
front-end development. The lower-latency live mode (\cref{sec:application})
uses an OpenCV display window driven by the same inference code, because
it avoids browser round-trips; it is a development tool rather than a second
interface framework.

\subsection{Technology stack}
\Cref{tab:stack} lists each technology, its purpose and where it is used.

\begin{table}[H]
\centering
\caption{Technology stack \rec.}
\label{tab:stack}
\small
\begin{tabularx}{\textwidth}{@{}p{3.3cm}Xp{3.9cm}@{}}
\toprule
\textbf{Technology} & \textbf{Purpose} & \textbf{System component} \\
\midrule
Python 3.11 & Implementation language for all modules & Entire system \\
PyTorch, torchaudio, torchvision & Model definition, training, CTC loss, audio
features, pretrained EfficientNet-B0 & Encoders, fusion, training, inference \\
OpenCV (\texttt{opencv-python}) & Video decoding, camera capture, cropping,
resizing, live display & Visual pipeline, live demo \\
MediaPipe & Face Mesh landmark detection~\cite{lugaresi2019,kartynnik2019} &
Mouth localisation \\
NumPy & Array manipulation and on-disk storage of cached crops & Preprocessing,
data loading \\
pandas & Dataset manifest (one row per utterance), split tables, result tables &
Data management, evaluation \\
soundfile & Reading and writing WAV audio & Audio pipeline, noise mixing \\
librosa & Reference implementation used to cross-check feature extraction in
tests & Unit tests \\
sounddevice & Microphone capture for the live loop & Live demo \\
scikit-learn & Speaker-grouped splitting (\texttt{GroupShuffleSplit}), confusion
matrices for extensions & Data splits, extensions \\
jiwer & Word and character error rate computation & Evaluation \\
Matplotlib & Training curves, WER-versus-SNR plots, attention-gate plots &
Experiment analysis \\
TensorBoard & Live logging of losses, error rates and learning rate & Training \\
PyYAML & Experiment configuration files & All scripts \\
pytest & Unit, integration and end-to-end tests & Quality assurance \\
Gradio & Web demonstration interface & Application \\
Git and GitHub & Version control and collaboration & Project management \\
\bottomrule
\end{tabularx}
\end{table}

\subsection{Dependencies and compatibility}
Three compatibility issues need attention when the environment is first
created.
\begin{itemize}
  \item \textbf{MediaPipe and Python versions.} MediaPipe publishes wheels for a
        limited range of Python versions. The team will confirm that the
        chosen Python version has a MediaPipe wheel for every member's
        operating system before fixing the version; Python~3.11 is proposed
        because it is widely supported \tbd.
  \item \textbf{PyTorch and CUDA.} The PyTorch build must match the CUDA driver
        of whichever GPU is used. The CPU-only build is sufficient for
        preprocessing, tests and the demonstration.
  \item \textbf{Paired library versions.} \texttt{torch}, \texttt{torchaudio}
        and \texttt{torchvision} must be installed as a matching set.
\end{itemize}

\subsection{Hardware requirements}
The project has not yet established which hardware is available. The plan
therefore states requirements rather than assuming a particular machine.
\begin{itemize}
  \item \textbf{Training} needs a CUDA-capable GPU; training the full model on a
        CPU alone is unlikely to finish within the six-week schedule. Candidate
        sources are a university GPU workstation, a team member's GPU, or a
        cloud notebook service \tbd. With a GPU of about 8\,GB memory, the
        recommended batch size of 16 clips with mixed precision is expected to
        fit; this will be confirmed by a memory test in week~2.
  \item \textbf{Storage.} Cached mouth crops for the whole corpus occupy roughly
        23\,GB as uncompressed 8-bit arrays
        ($\approx 33{,}000 \times 75 \times 96 \times 96$ bytes), plus the raw
        corpus and MUSAN. About 100\,GB of free disk space is recommended.
  \item \textbf{Demonstration} requires a laptop with a webcam and microphone. A
        GPU is helpful but not required; latency will be measured on the actual
        demonstration machine.
\end{itemize}

\subsection{Reproducibility measures}
\begin{itemize}
  \item Exact dependency versions are pinned in \texttt{requirements.txt} (and an
        equivalent \texttt{environment.yml}), generated from a working
        environment rather than written by hand.
  \item Every experiment is defined by a YAML configuration file committed to
        Git; the configuration and Git commit hash are saved inside each
        checkpoint.
  \item Random seeds for Python, NumPy and PyTorch are set from the
        configuration, and the noise-mixing seed for the test set is fixed so
        that every model is evaluated on identical noisy audio.
  \item The speaker split is stored as a file (\texttt{splits.csv}) and loaded by
        every script, never regenerated.
  \item Preprocessing scripts are deterministic and record the MediaPipe
        version and crop parameters in the dataset manifest.
\end{itemize}
